\documentclass[11pt,a4paper]{letter}

\usepackage[utf8]{inputenc}
\usepackage[T1]{fontenc}
\usepackage[english]{babel}
\usepackage{geometry}
\usepackage{xcolor}
\usepackage[pdfborder={0 0 0}]{hyperref}
	\newcommand{\qm}[1]{“#1”}


\usepackage{fontspec}
\setmainfont{Times New Roman}
\newfontfamily\vest{TRY Vesterbro Regular}
\newfontfamily\vestb{TRY Vesterbro Poster}
\newfontfamily\overp{Overpass}
\geometry{margin=2.5cm}
\definecolor{letterblue}{RGB}{30, 89, 111}
\let\letterclassopening\opening
\renewcommand{\opening}[1]{\letterclassopening{#1}\vspace{2em}}
\renewcommand{\closing}[1]{\par\vspace{2em}\noindent #1\par\vspace{1em}\noindent\fromsig\par}

\signature{Anders Thaysen}
\address{
{\color{letterblue}\vestb Anders Thaysen}
}
\date{\vest\today}

\begin{document}

\begin{letter}{}
\opening{Kære ledelse på Næstved Gymnasium,}

Denne klage giver for mig ikke anledning til at genoverveje vores daværende vurdering af standpunkt.

Det er korrekt at min computer løb tør for strøm undervejs i eksaminationen. Jeg kan dog ikke vedkende mig at det skulle have \qm{påvirket} forløbet på nogen måde. Hændelsen varede ca. 7 sekunder, og mit fokus var uændret undervejs og efterfølgende.

Jeg er ikke enig i elevens beskrivelse af dialogen i eksaminationsforløbet. Den manglende forståelse bundede efter min overbevisning i en faglig usikkerhed hos eleven. Engelskfaglige spørgsmål om f.eks. afsenderforhold i forbindelse med fremstilling af Lance Armstrong i analyseobjektet blev ikke forstået. Ej heller videnskabsteoretiske spørgsmål om de engelskfaglige metoders rolle i besvarelsen af opgavens problemformulering blev forstået.

Det er korrekt at det var vores vurdering at analysen i opgaven ikke var tilstrækkeligt \qm{dybdegående}, og at elevens oplæg ikke ændrede ved denne vurdering. Analysen af dokumentaren var f.eks. uden engelskfaglige begreber, og dens konklusioner var vage og oftest helt uden belæg. 

Opgaven har dog, som der også står i klagen, en slags besvarelse af \qm{alle dele af opgaveformuleringen}, og jeg fastholder derfor, at karakteren 02 er den rette karakter for den pågældende elev.

\closing{Venlig hilsen,}

\end{letter}

\end{document}
